% =====================================================================
% Stage 5: Title options, Abstract, Limitations, Conclusion
% =====================================================================

% ---------------------------------------------------------------------
% TITLE OPTIONS (choose one)
% 1. Same Planet, Three Codes: Inter-Code Systematics in Exoplanet
%    Transmission Spectra and Their Effect on Machine-Learning Retrievals
% 2. How Much Does the Forward Model Matter? A Paired Comparison of
%    POSEIDON, petitRADTRANS, and TauREx over 98,725 Simulated Planets
% 3. Forward-Model Mismatch as Simulator Misspecification: Detectability
%    and Retrieval Bias across Three Radiative-Transfer Codes
% 4. Detecting and Quantifying Radiative-Transfer Code Systematics in
%    Ariel-Like Transmission Spectra
% 5. Low-Dimensional, Deterministic, and Persistent: Inter-Code
%    Systematics in Simulated Exoplanet Transmission Spectra
% ---------------------------------------------------------------------

% The abstract lives in abstract.tex so it can precede the Introduction.


% =====================================================================
\section{Limitations}
\label{sec:limitations}

\paragraph{Simplified atmospheres.}
All spectra assume an isothermal, cloud-free atmosphere, and the codes use their default opacity databases.
Clouds, hazes, non-isothermal profiles, and disequilibrium chemistry would introduce further modelling choices on which the codes may disagree.
Inter-code differences for realistic atmospheres may therefore be larger than those reported here.
The CH$_4$ and NH$_3$ abundances are restricted above 2000\,K and 1600\,K because opacity data are unavailable there, which leaves part of the hot regime sparsely covered for those species [TODO: VERIFY thresholds].

\paragraph{Confounded sources of disagreement.}
The three codes differ in opacity data, in opacity treatment (correlated-$k$ at $R=1000$ for petitRADTRANS, cross-sections at $R=15{,}000$ for the other two), and in internal conventions.
Our paired design measures their combined effect.
It does not attribute a residual to a single cause, and the physical interpretations in Section~\ref{sec:res_q1} are hypotheses that require controlled re-simulation to test.

\paragraph{Noise model.}
Observational noise is white, Gaussian, and uniform across bins.
Real Ariel and JWST errors vary with wavelength and include correlated instrumental and stellar systematics, which can either mask or mimic inter-code differences.

\paragraph{Sample selection.}
The $210\,\mathrm{ppm}$ feature-height cut retains $44.6\%$ of planets and preferentially removes cool planets (18.6\% retained below 600\,K, 55.1\% above 2000\,K).
Population-level statistics therefore describe large-feature planets, and the hardest planets for the autoencoders are under-represented in training.

\paragraph{Model and statistical scope.}
Retrievals are point estimates from gradient-boosted trees trained without noise, not Bayesian posteriors.
Most of these regressors reached the cap of 3000 boosting rounds, so they are not fully converged.
The random-forest classifier fits its training set perfectly.
Results use a single train/validation/test split, two to three initialisation seeds per autoencoder configuration, and one seed for the TauREx detector.
The width--schedule ablation varies the learning-rate schedule and the early-stopping patience jointly.
Ten of the twelve denoising models stopped early, so the negative result for the denoising objective may partly reflect the stopping rule rather than the objective.

% =====================================================================
\section{Conclusion}
\label{sec:conclusion}

We measured the systematic differences between three widely used radiative-transfer codes on a paired sample of $98{,}725$ simulated planets, and traced their consequences for detection and for machine-learning retrieval.
The differences consist of a depth offset of about $20\,\mathrm{ppm}$ between POSEIDON and the other two codes, which scales with feature height, and a shape residual of about $1\%$ of the feature height.
The shape residual is concentrated in a few spectral regions and is a near-deterministic function of the planet parameters.
Without noise these differences are easily detected, even when they are only a few ppm.
At realistic noise levels they become difficult to detect from a single spectrum, and the difficulty is set by the detector: one-class autoencoders are limited by their own worst-case reconstruction errors, whereas a classifier trained on all three codes approaches or exceeds a parameter-aware oracle.
In retrieval, code mismatch produces parameter-dependent biases that additional training data from a single code do not remove.
For individual noisy spectra these biases are small compared with the statistical scatter, but for population studies they persist as coherent shifts.

These results suggest four practices.
First, ML retrievals should be validated on spectra from forward models other than the one used for training, and should report the resulting self/cross gap.
Second, because the residuals are low-dimensional and predictable from the parameters, they could be emulated and either corrected or included as a modelled error term; we have not tested this here.
Third, code identification should use discriminative models trained on all candidate codes rather than one-class reconstruction scores.
Fourth, reconstruction-based detectors should be evaluated at a calibrated false-positive rate, because their AUC can greatly overstate their usable detection rate.
Extending the paired design to cloudy, non-isothermal atmospheres, shared opacity sources, and instrument-specific noise would separate the physical origins of the systematics and test whether these conclusions hold for real observations.
